\begin{donnees}
    \begin{itemize}
        \item $m=\qty{255}{\gram}$~;
              Tous les frottements sont négligeables.
    \end{itemize}
\end{donnees}
On écarte $\mathrm{(S)}$ (Fig.~\ref{fig:solide_ressort}) de sa position
d'équilibre d'une distance $X_m$ et on l'abandonne sans vitesse initiale.
L'équation horaire du mouvement de $G$ est~:
${x(t)=X_m \cdot \cos\left(\frac{2\pi}{T_0} \cdot t+\varphi\right)}$.
\begin{enumerate}
    \item L'équation de la vitesse de $G$ s'écrit~:
          $v(t)=-0,25 \cdot \sin(2\pi\cdot t)$ \; (\unit{\v}).
          \begin{enumerate}
              \item En exploitant l'équation de la vitesse, déterminer la
                    période propre $T_0$ des oscillations, la valeur de
                    l'amplitude $X_m$ et la phase $\varphi$ à $t_0=0$.
              \item Vérifier que la raideur du ressort est
                    $K\approx\qty{10}{\newton\per\meter}$.
          \end{enumerate}
    \item Déterminer l'expression de la force de rappel $\vec{F}$ exercée par le
          ressort sur le solide $\mathrm{(S)}$ à l'instant
          $t=\qty{0,5}{\second}$.
\end{enumerate}

